\section{Descripción del problema}
\label{sec:descripcion-problema}

% Sección agregada a la plantilla institucional. Aquí se plantea de forma
% precisa y acotada el problema específico que la tesis atacará, a partir
% de lo expuesto en la Introducción. Se recomienda responder:
%   - ¿Cuál es el problema exacto (no el área general)?
%   - ¿Quién lo sufre / en qué contexto se manifiesta?
%   - ¿Qué evidencia hay de que es un problema real y no resuelto?
%   - ¿Qué pasa si no se resuelve (consecuencias)?

En la actualidad, la ingeniería de software está migrando hacia lo que, de acuerdo con lo que retoma \textcite{fernandezyfernandez2025rsm} de Hassan et al., se denomina SE 3.0 o "ingeniería de software agéntica": un enfoque en el que el desarrollador deja de escribir código directamente y pasa a actuar como director de agentes autónomos. Dentro de este contexto surge el llamado \textit{vibe coding}, término acuñado en febrero de 2025 por Andrej Karpathy, cofundador de OpenAI, para referirse a la posibilidad de "olvidar que el código existe" \parencite{bbc2026vibecoding}. Según \textcite{fernandezyfernandez2025rsm}, esta práctica describe una forma de programación conversacional en la que el desarrollador deja que el modelo tome el control de la implementación, priorizando el resultado obtenido sobre la comprensión del código producido, lo cual ha generado dudas sobre la confiabilidad arquitectónica del código resultante y sobre un posible deterioro, a largo plazo, de las competencias técnicas de los desarrolladores.

Como respuesta a estos retos han surgido distintas estrategias de coordinación multiagente dentro de la ingeniería de software agéntica. Sistemas como MetaGPT \parencite{hong2024metagpt} y ChatDev \parencite{qian2024chatdev} plantean una organización basada en roles, donde los agentes especializados participan en diferentes etapas del desarrollo. MetaGPT utiliza procedimientos operativos estandarizados mediante secuencias de prompts para disminuir la propagación de errores que puede surgir al encadenar modelos de forma ingenua, mientras que ChatDev organiza la comunicación entre agentes especializados en las fases de diseño, codificación y pruebas mediante una estructura de diálogo. Siguiendo este mismo principio, \textcite{fernandezyfernandez2025rsm} proponen el Role Specialization Model (RSM), que lleva la especialización de roles hacia las propias herramientas de asistencia: cada herramienta recibe una responsabilidad específica (arquitecto, analista, especialista) de acuerdo con sus capacidades, de modo que sus contribuciones se complementen y se minimice el solapamiento entre funciones.

Sin embargo, esta separación de responsabilidades no se sostuvo por completo en la ejecución real del RSM. \textcite{fernandezyfernandez2025rsm} identifican, entre otros factores, el costo cognitivo asociado con cambiar de herramienta a mitad de una tarea, que llevó a que el rol de analista asumiera responsabilidades propias del especialista, como la refactorización arquitectónica. Por otra parte, \textcite{feng2025heteroswarms} muestran que la selección conjunta del modelo asignado a cada rol dentro de un sistema multiagente puede mejorar de manera considerable el desempeño frente a configuraciones homogéneas. No obstante, en este último caso, la heterogeneidad se plantea a nivel de los modelos utilizados y no de las herramientas ni de sus capacidades para acceder a recursos externos.

A partir de la literatura revisada, se observa que tanto la especialización de roles como la heterogeneidad de modelos por rol han sido estudiadas de manera independiente. Sin embargo, la combinación de roles especializados con toolsets heterogéneos, es decir, con conjuntos de herramientas diferentes asignados de acuerdo con la responsabilidad de cada rol, no ha sido estudiada de forma aislada como variable experimental dentro de un entorno real de desarrollo de software.


\subsection{Planteamiento del problema}
% Enunciado claro y acotado del problema.

\subsection{Evidencia del problema}
% Datos, casos, referencias que sustentan que el problema existe.

\subsection{Consecuencias de no resolverlo}
% Impacto de dejar el problema sin atender.
